\ifdefined\Needspace\Needspace{25\baselineskip}\fi
\begin{samepage}
\begin{exercisechunk}
\item \textbf{The retention cost of a compulsory probability floor.}
\label[exercise]{ex:l21-probability-floor}
\exerciseprofile{2}{2}{2}{\exproof\quad\excalculation\quad\exinterpretation}{%
Compare retention between two learners with the same order of continuation regret.}
In the two-action full-information problem, maintain $Q_0=1/2$ and
update $Q_t$ by the projected-gradient rule \cref{eq:projected}, with
$Q_t$ in place of $X_t$ and $\eta=T^{-1/2}$.
For a fixed $0<\alpha\le1$, play action one with probability
$X_t=(1-\alpha)Q_t+\alpha/2$; the complete loss vector is still observed.
The update of $Q_t$ uses that vector and is independent of the sampled action.
Call the resulting algorithm $\cA_{T,\alpha}$.
\begin{enumerate}[(a)]
\item Prove
$\operatorname{Cont}_{\cP_{\rm fi}}^{\cA_{T,\alpha}}(T)\le\sqrt T+\alpha T/2$.
\item On the constant loss vector $c=(0,1)^\top$, prove
$\operatorname{Ret}_{\delta_c}^{\cA_{T,\alpha},\infty}(K)=\alpha K/2$.
\item Set $\alpha=1/T$. Compare the continuation order and tail retention
with those of the unmodified projected-gradient learner.
Explain why this comparison rules out the claim that a positive probability
floor is necessary for optimal-order continuation under full feedback.
\end{enumerate}
\begin{hints}
Compare the mixture losses at $X_t$ and $Q_t$ before using the projected
gradient bound.
\end{hints}
\end{exercisechunk}
\end{samepage}
